Approximately 600 million people in sub-Saharan Africa lack electricity access. Solar mini-grids offer practical solutions, but financing remains a major barrier. Carbon markets have been proposed as a way to generate additional revenue by selling carbon credits from emission reductions. However, little evidence exists on whether this approach actually works for small-scale rural energy projects. This study examines two questions: How effectively does carbon finance improve mini-grid investment and commercial viability, and what conditions are necessary for successful integration at scale? Through nine interviews with developers, carbon experts, policy officials, and registry professionals, plus focus groups with 30 community members across three Nigerian sites, the research reveals that carbon finance is supplementary rather than transformative. Developers did not base their investment decisions on carbon revenue. For projects under 200 kW, certification costs represent 8\% of capital expenditure, while current prices of \$1-5 per credit mean transaction costs exceed potential revenues. The atmosfair–MiniGridCo programme shows that combining multiple mini-grids into a single carbon project is technically possible when institutional support is available. However, the fact that it has not been replicated across more than 120 other Nigerian mini-grids points to a market failure. The conditions that made it work are not accessible to most developers. Successful integration requires coordinated public intervention across five areas: technical capacity, financial viability, regulatory frameworks, institutional arrangements, and carbon market design. Without structural transformation including publicly-funded aggregation and regulatory reforms, carbon markets remain largely inaccessible to mini-grids.